%! Setting Up a Test for the Difference of Two Population Proportions — Unit 6, Lesson 6.10 — Guided Notes
\documentclass[10pt]{article}
\usepackage{apstats-article}
\usepackage{apstats-boxes}
\newcommand{\TallMath}[1]{$\displaystyle #1\rule[-1.4em]{0pt}{3.2em}$}

\begin{document}

\pageheader{Unit 6, Lesson 6.10}{Guided Notes}
\namedateperiod

\vspace{0.15in}

\begin{objectivebox}
\textbf{Lesson 6.10 Objectives:} Set up a two-sample $z$-test for the difference of two population
proportions: state hypotheses, compute $\hat p_c$, and verify conditions. \hfill \textit{VAR-6.H--J, Skills 1.E \& 4.C}
\end{objectivebox}

\begin{vocabbox}
\begin{description}
  \item[\termblanklong{pooled proportion ($\hat p_c$)}]
        The combined success proportion under $H_0: p_1=p_2$:
        $\hat p_c = \dfrac{n_1\hat p_1 + n_2\hat p_2}{n_1+n_2}$. Used for Large Counts in a two-sample test.
\end{description}
\end{vocabbox}

\begin{notesbox}{Hypotheses}
$H_0$: \blank{6cm} \quad ($p_1 - p_2 = $ \blank{1.5cm})\\[6pt]
$H_a$ (choose one): \quad
$p_1 \blank{1cm} p_2$ \quad or \quad $p_1 \blank{1cm} p_2$ \quad or \quad $p_1 \blank{1cm} p_2$
\end{notesbox}

\begin{notesbox}{Pooled Proportion}
Under $H_0: p_1=p_2$, both groups share one unknown proportion. We estimate it by pooling:
\[\hat p_c = \frac{n_1\hat p_1 + n_2\hat p_2}{n_1+n_2}
           = \frac{\text{\blank{4cm}}}{\text{\blank{3cm}}}\]
Use $\hat p_c$ (not $\hat p_1$ or $\hat p_2$) for the \textbf{Large Counts} condition.
\end{notesbox}

\begin{notesbox}{Conditions}
\renewcommand{\arraystretch}{1.7}
\begin{tabular}{@{} p{0.22\linewidth} p{0.72\linewidth} @{}}
\textbf{Random}     & \blank{8cm} \\
\textbf{Large Counts} & All four values $\ge 10$: \\
                    & $n_1\hat p_c = $ \blank{2.5cm}, \quad $n_1(1-\hat p_c) = $ \blank{2.5cm} \\
                    & $n_2\hat p_c = $ \blank{2.5cm}, \quad $n_2(1-\hat p_c) = $ \blank{2.5cm} \\
\textbf{10\%}       & $n_1 \le 10\%$ of pop.\ 1 \textbf{and} $n_2 \le 10\%$ of pop.\ 2 \\
\end{tabular}
\end{notesbox}

\begin{notesbox}{Worked Example --- App Error Rates}
Group A (new app): 42 errors in 300 tasks. Group B (old app): 60 errors in 300 tasks.
Test at $\alpha=0.05$: do the error rates differ?

\textit{State:} $H_0$: \blank{6cm}; $H_a$: \blank{6cm}

\textit{Plan:} $\hat p_A = $ \blank{2cm}, $\hat p_B = $ \blank{2cm}. \quad
$\hat p_c = \dfrac{\blank{3cm}}{\blank{2cm}} = $ \blank{2cm}

\textbf{Random:} \blank{6cm} \\
\textbf{Large Counts:} $300(\blank{2cm}) = $ \blank{2cm}, $300(1-\blank{2cm}) = $ \blank{2cm} \quad (same for both groups).\\
\textbf{10\%:} \blank{6cm}
\end{notesbox}

\begin{practicebox}
\textbf{Try It:} Researchers compare pass rates for a certification exam between two training
programs. Program 1: 56 of 80 passed. Program 2: 44 of 80 passed.
Is there evidence that the pass rates differ ($\alpha=0.05$)?

State $H_0$ and $H_a$, compute $\hat p_c$, and verify conditions.
\writelines{5}
\end{practicebox}

\end{document}
